\documentclass[11pt]{article}
\usepackage[margin=1in]{geometry}
\usepackage{amsmath}

\title{Heat Flow in a Thin Rod}
\author{A.~Student}
\date{\today}

\begin{document}
\maketitle

\begin{abstract}
We derive the one-dimensional heat equation
for a thin, insulated rod and solve it by
separation of variables. The solution is a
sum of decaying Fourier modes.
\end{abstract}

\section{The heat equation}
Let $u(x,t)$ be the temperature of a rod of
length $L$ with diffusivity $\alpha > 0$.
Conservation of energy and Fourier's law give
\begin{equation}
  \frac{\partial u}{\partial t}
    = \alpha \, \frac{\partial^2 u}{\partial x^2},
  \qquad 0 < x < L.
  \label{eq:heat}
\end{equation}
With the ends held at zero, \eqref{eq:heat} has the solution
\begin{equation}
  u(x,t) = \sum_{n=1}^{\infty} b_n
    \sin\!\left(\frac{n\pi x}{L}\right)
    e^{-\alpha (n\pi/L)^2 t}.
\end{equation}

\section{Separation of variables}
Substituting $u(x,t) = X(x)\,T(t)$ into \eqref{eq:heat} and dividing
by $\alpha X T$ separates the equation into two ordinary differential
equations, one in space and one in time:
\begin{equation}
  X'' + \lambda X = 0, \qquad T' + \alpha \lambda T = 0.
\end{equation}
The boundary conditions $X(0) = X(L) = 0$ admit nonzero solutions only
when $\lambda_n = (n\pi/L)^2$ for $n = 1, 2, \dots$, giving the
eigenfunctions $X_n(x) = \sin(n\pi x / L)$. Each time factor decays as
$T_n(t) = e^{-\alpha \lambda_n t}$, so higher modes vanish fastest and
the temperature profile smooths out over time.

The coefficients $b_n$ follow from the initial temperature $f(x) = u(x,0)$
by orthogonality of the sine functions on $[0, L]$:
\begin{equation}
  b_n = \frac{2}{L} \int_0^L f(x) \sin\!\left(\frac{n\pi x}{L}\right) dx.
\end{equation}

\section{Discussion}
Because every mode decays exponentially, the rod approaches the uniform
temperature of its ends. The slowest mode, $n = 1$, sets the time scale
$\tau = L^2 / (\alpha \pi^2)$ on which this happens.

\end{document}
% Rebuild ../sample-page-1.png from this directory with
% pdflatex paper.tex && pdflatex paper.tex && pdftoppm -png -singlefile -scale-to-x 1020 -scale-to-y -1 paper.pdf ../sample-page-1
